\documentclass[a4paper,11pt]{article}

\usepackage{revy}
\usepackage[utf8]{inputenc}
\usepackage[T1]{fontenc}
\usepackage[danish]{babel}
\usepackage{amsmath,amssymb}

\revyname{MatRevy}
\revyyear{2026}
\version{2026-10-09}
\eta{4 minutter}

\title{Hilberts Hotel 2: Rengøring}
\author{Julius Pallesen'24}

\begin{document}
\maketitle

\begin{sketch}
\scene{Scenen er Hilberts Hotel. Der er tre døre med numre 1,2,3 på scenen og et skilt til værelser $\mathbb N\cap[4,\infty)$ der peger ud af scenen bag sidetæppet. O står med sin rengøringsvogn. H viser M hen til O}

\says{H} Lad mig introducere dig for O, hotellets oldfrue. O har gjort rent dagligt på samtlige værelser siden tidernes morgen, da dette hoteleventyr begyndte. Nu vil hun give stafetten videre, og lære dig hvordan man gør.

\scene{H skynder sig ud af scenen, da han er en travl mand.}

\says{M} \act{Kalder efter H der er på vej væk} Men, men.. Jeg troede ikke jeg skulle gøre rent. Jeg troede bare jeg skulle fortælle gæsterne hvilke værelser de skulle flytte til. Jeg har en kandidatgrad i matematik!

\says{O} Det har vi sgu da allesammen. Det er der ikke noget imponerende i. Grib så en moppe. Vi har meget vi skal nå.

\says{M} Altså.. Jeg er totalt Sisyfos. Nu skal jeg gøre uendelig mange værelser rent til evig tid :(

\says{O} Ti så stille. Hvis man gør det ordentligt tager det kun to minutter.

\says{M} To minutter? Man kan da ikke gøre uendelig mange værelser rent på to minutter.

\says{O} Det kan man i hvert fald. Kig og lær.

\scene{O går hen til dør 1 og ringer på.}

\says{O} Jeg gør rent herinde i ét minut.

\scene{O går ind og vi høre stemmerne af O og G1.}

\says{G1} Hey jeg er lige flyttet ind. Kan I gøre rent senere?

\says{O} Jeg har et travlt program, du bliver nødt til lige at flytte dig lidt.

\says{G1} Men jeg har ikke engang pakket ud!?

\says{O} Spar dig besværet. Du kommer til at skulle rykke værelse lige om lidt.

\scene{H kommer tilbage på scenen med G0, tager ham ind på værelse 1}

\says{H} Du skal bo her.

\says{M} Chef! Jeg tænkte bare...

\says{H} Det skal du holde op med! Bare gør hvad O siger. Jeg skal lige rykke rundt på nogle gæster. \act{Skifter til at være høflig overfor G1} Jeg beklager virkelig, men jeg har desværre vist dig hen til det forkerte værelse. Hvis du lige vil komme med mig.

\scene{Resten af værelsesrokaden finder sted i stilhed. Imens er ét minut er gået (det behøver ikke være et rigtigt minut), og O går ud og videre til dør 2 og ringer på.}

\says{O} Så gør jeg rent herinde i et halvt minut.

\scene{M prøver at følge med}

\says{O} Du skal nok få lov. Nu kigger du bare, og så lader du mig arbejde.

\says{M} Men hvorfor skulle jeg så tage en moppe?

\scene{O ignorerer og går ind på værelse 2.}

\says{G1} Altså hvad er meningen med det her?

\says{O} Du vil måske hellere have lus og fnat, end sæbe og håndklæder? Nej det tænkte jeg nok!

\scene{O går videre til værelse 3. Halvandet minut er nu gået.}

\says{O} Og så går jeg herind i et kvartminut.

\says{G2} \act{Skriger}

\says{O} Tror du måske aldrig, jeg har set sådan en nøgen mand/dame før?

\scene{Der er nu gået 1:45. O kommer ud af værelset}

\says{O} Prik prik prik

\scene{O går ud ad side tæppet til de resterende værelser. Der lyder en klokke hvert $1/2^n$'te minut. Når to minutter er gået, kommer O tilbage}

\says{O} Se det var da let nok. Så er det din tur!

\says{M} Jamen nu har du jo lige taget alle værelserne

\says{O} Nej, jeg efterlod alle lige primtal over 2 til dig. Ja kom så i gang!

\scene{M, forvirret, skynder sig mod de høje tal}

\says{O} Hahaha hvor længe mon der går før han finder ud af det.

\scene{Gk kommer ind og er vred}

\says{Gk} Hey! Er det dig der er rengøringsdamen her på hotellet?

\says{O} Jeg er oldfrue! Og hvem er du?

\says{Gk} Jeg bor på værelse 43 trillioner 252 billiarder 333.

\says{O} Gud! Det er et af de små tal!

\says{Gk} Se her!

\scene{Gk viser sin finger som om der ligger noget på den}

\says{O} Hvad er det?

\says{Gk} JA HVAD ER DET! Der hvor jeg kommer fra ville man kalde det en nullermand. Sig mig gør I ikke rent ude i hjørnerne?

\scene{M kommer tilbage på scenen}

\says{M} Jeg kunne altså ikke finde nogen lige primtal udover 2

\scene{Hverken O eller Gk kan holde masken af grin, men prøver hårdt}

\says{O} Det er lige meget nu. Alle værelserne trænger alligevel til en ekstra omgang. Kan du huske hvordan man gjorde?

\says{M} Det tror jeg nok. Så det var ét minut på værelse 1, et halvt minut på værelse 2, en tredjedel af et minut på værelse 3.. Prik prik prik?

\scene{O og Gk kan igen næsten ikke holde masken, men formår at skjule det for M}

\says{O} Ja præcis. Du går bare i gang. Husk at nå ud i hjørnerne!

\scene{M er går ind på værelse 1}

\says{O\&Gk} HAHAHAHAHAHA!!!!

\says{Gk} Du er så ond! Lige primtal? Haha

\says{O} Og harmonisk rengøring! Hahaha han bliver aldrig færdig!

\says{Gk} Han er totalt Sisyfos

\scene{Lys ned. Fortsættes...}
\end{sketch}
\end{document}
